\section{Multiple Comparisons and Bayesian Power Analysis}\label{app:multiple_comparisons}

\subsection*{Multiple Comparisons}

A practical advantage of the HDI+ROPE framework over frequentist approaches concerns
multiple comparisons. In NHST, conducting several tests simultaneously inflates the
Type~I error rate, requiring corrections (e.g.\ Bonferroni) that reduce the power of
each individual test. Frequentist confidence intervals used for simultaneous inference inherit the same
dependence on the full set of intended comparisons, since their simultaneous coverage
properties are tied to p-values. Consequently, both power and precision planning must account for the number
of planned tests.

The Bayesian posterior is not affected in this way. Because the posterior distribution
is determined solely by the observed data and the assumed data-generating process ---
not by the set of tests that were intended but not performed --- the power of any one
test is unaffected by how many other comparisons are planned, provided each test uses
an independent model \citep{gelman2012, kruschke2015doing}.
An independent model means each comparison is fitted separately, without shared
parameters or a common prior; when tests are connected through a hierarchical model,
posteriors are linked across comparisons and the independence no longer holds
\citep{gelman2012}.
% \citet{gelman2012}: Gelman, Hill \& Yajima (2012) is the standard journal reference
% for the claim that Bayesian posteriors typically require no multiple-comparison
% correction when tests are based on independent models, and also covers the
% hierarchical case where posteriors are linked and the argument requires modification.
% \citet{kruschke2015doing}: DBDA2 (Chapter~11) provides the pedagogical treatment.
The precision goal $\omega_{\rm goal}$ and the stopping rule of DPitG therefore apply
unchanged in a multiple-testing context, without the need for correction.

\subsection*{Three Types of Bayesian Power Analysis}

The same design flexibility extends to how power is conceptualised in the Bayesian
setting, where \citet{kruschke2015doing} distinguishes three types of power analysis that arise in
the Bayesian setting.

\textit{Prospective} power analysis is conducted before data collection: given a
hypothesised data-generating distribution (from theory, pilot data, or prior research),
one simulates experiments and asks how often the precision goal $\omega_{\rm goal}$ is
achieved within the budget $N_{\rm max}$.
The analytical formula for $N_{\rm goal}$ (Equation~\ref{eq:pitg_stop_iteration})
provides this as a closed-form lower bound on the required sample size; DPitG may
require additional samples when the true parameter value lies near a ROPE boundary
(see Section~\ref{sec:trends}).

\textit{Retrospective} power analysis uses the posterior from already-collected data as
the hypothetical data generator to ask whether the completed study was sufficiently
powered. In NHST this quantity is a monotone function of the p-value and therefore adds
no new information beyond it \citep{hoenig2001}; in the Bayesian setting it can reveal
whether the precision goal $\omega_{\rm goal}$ was realistically achievable given the
observed data, independently of the decisiveness verdict.
% \citet{hoenig2001}: Hoenig \& Heisey (2001) demonstrate that observed power in NHST
% is a 1-to-1 function of the p-value and therefore adds no information beyond it ---
% the canonical reference for the redundancy of retrospective NHST power.

\textit{Replication} power asks: if the experiment were repeated exactly, what is the
probability of achieving the goal in the new study, informed by the results of the
first? In the DPitG framework this is straightforward to compute by using the initial
posterior as the prior for the follow-up and re-running the stopping rule; NHST has
no principled access to a posterior-derived data generator and cannot support this
form of analysis \citep{kruschke2015doing}.
% \citet{kruschke2015doing}: DBDA2 (Chapter~13) introduces replication power as the
% probability of achieving the Bayesian goal in a follow-up study seeded by the
% current posterior --- the direct justification for citing Kruschke here rather than
% a frequentist replication reference.
